\documentclass{article}

\makeatletter
\def\input@path{{Formatting_Instructions_For_NeurIPS_2026/}}
\makeatother
\usepackage[dblblindworkshop]{neurips_2026}
\usepackage[utf8]{inputenc}
\usepackage[T1]{fontenc}
\usepackage{hyperref}
\usepackage{url}
\usepackage{booktabs}
\usepackage{amsmath}
\usepackage{microtype}
\usepackage{graphicx}

\title{Learning Energy-Harvesting Soaring Policies in Time-Varying Turbulent Wind}
\workshoptitle{RL for Experimental Sciences Workshop}

\author{Anonymous Authors}

\begin{document}

\maketitle

\begin{abstract}
Soaring converts spatial and temporal variation in the atmosphere into usable
energy, but learning such behavior is difficult when the useful signal is local,
noisy, and coupled to vehicle dynamics.  We introduce ThermalHunter, a
controlled numerical platform for studying energy-harvesting soaring in
time-varying turbulent wind.  The platform couples DNS snapshots of
Rayleigh--Benard convection at $Ra=5\times10^{7}$ to two complementary flight
  models: a transparent quasi-steady discrete controller and a higher-fidelity
  dynamic aerodynamic environment.  In an exploratory steady-regime comparison
  that reused the model-selection scenarios, a tuned DQN achieved a mean height
  change of $337.4$ m (bootstrap 95\% CI $[297.4,376.2]$), compared with
  $-107.1$ m for random control.  We also show that DQN action maps can suggest
  sensor thresholds for an interpretable tabular Q-learning controller; that
  derived controller attained $192.3$ m on the same scenarios.  In the
  higher-fidelity dynamic regime, three PPO
  replicas evaluated on a new 100-scenario suite had mean height changes of
  $397.7$, $276.8$, and $378.5$ m, while three dynamic-DQN replicas had
  $-146.3$, $-121.6$, and $-76.0$ m.  These results support energy-harvesting
  behavior in the simulated turbulent flow, but not route navigation, safety,
  real-flight deployment, or sim-to-real transfer.
\end{abstract}

\section{Introduction}

Unpowered gliders can gain energy by exploiting atmospheric motion.  The
problem is attractive for reinforcement learning (RL) because the controller
must repeatedly choose how to trade local lift, turning, and aerodynamic loss
without an externally supplied trajectory.  Field work has shown that an RL
controller can support autonomous glider soaring \citep{reddy2018}, while
simulation studies have demonstrated that learned policies can control gliding
trajectories with nontrivial aerodynamics \citep{novati2019}.  Those results
motivate a more controlled question: can an agent learn to harvest energy from
time-varying, turbulent flow using only quantities that could plausibly be
available to an onboard controller?

ThermalHunter addresses that question with a numerical experimental platform.
Its wind input consists of direct numerical simulation (DNS) snapshots of
Rayleigh--Benard convection at $Ra=5\times10^{7}$.  The wind is therefore
simulation data, not an atmospheric measurement or a model of a particular
flight test.  The platform exposes local vertical-wind acceleration and a
wing-span vertical-wind difference, rather than the complete flow field, and
uses height change and total-energy-height change as outcomes.  It is designed
to study energy extraction from a turbulent environment, not to optimize a
route or demonstrate a safety-critical autonomous aircraft.

The study has two complementary fidelity regimes.  The first combines a
quasi-steady glider model with discrete angle-of-attack and bank commands.  It
provides a small, inspectable testbed for comparing tabular Q-learning and DQN.
The second retains dynamic flight and actuator states and tests whether the
learned energy-harvesting behavior persists under a richer model.
The regimes are not algorithmic leaderboards: policies are compared only within
the environment in which they were trained and evaluated.

This manuscript makes three bounded contributions.  First, it specifies a
two-fidelity RL testbed driven by a single DNS wind source and matched-scenario
evaluation.  Second, it establishes a steady-regime baseline in which the DQN
is evaluated against random and tabular-Q controllers under identical initial
conditions.  Third, it treats the DQN action map as a representation-design
tool: its coarse decision boundaries are used to choose transparent sensor bins
for a tabular controller.  The last comparison is reported as preliminary,
because it has not yet been replicated across independent training seeds or an
untouched test suite.

\section{Related work}

RL has been applied to flight control and to decision-making in complex flows.
Reddy et al.\ \citep{reddy2018} demonstrated glider soaring via RL in the
field, linking learned control to an operational aircraft setting.  Novati et
al.\ \citep{novati2019} used deep RL for controlled gliding and perching,
showing how learning can accommodate aerodynamic control objectives.  In a
different flow-physics setting, Colabrese et al.\ \citep{colabrese2017}
showed that RL can discover navigation strategies from local observations in
vortical flows.  ThermalHunter shares the local-observation premise but changes
the objective: the primary outcome is energy or altitude gain, not displacement
toward a target.

The present work also uses a standard value-based deep RL baseline.  DQN
combines experience replay and a target network to learn action values from
high-dimensional observations \citep{mnih2015}.  Here its role is deliberately
modest.  The DQN is not presented as a universal control solution; instead, it
offers a continuous sensor-state policy against which a compact discretization
can be inspected.  That comparison exposes a practical experimental question:
whether a learned action map can help select a lower-dimensional, auditable
state representation for tabular learning.

\section{ThermalHunter environments}
\label{sec:environment}

\subsection{Wind input and scaling}

The shared input is a sequence of three-dimensional velocity fields from DNS
Rayleigh--Benard convection at $Ra=5\times10^{7}$.  At each query, the simulator
uses trilinear spatial interpolation; the dynamic environment additionally
interpolates between adjacent temporal frames.  Horizontal position is periodic
in the $1000\times1000\times1000$ m numerical domain.  To give the flow a
consistent relation to the vehicle, the simulator scales the three-dimensional
wind so that its RMS speed is $0.5$ times the trim true airspeed at the median
angle of attack.  This scaling is a modeling choice, not a claim about an
observed atmosphere.

Each rollout begins from a randomized horizontal position, altitude, and
heading.  During the steady-regime evaluation, the initial wind frame is drawn
from frames 60--99, excluding the first 60 DNS frames.  The initial altitude is
drawn uniformly from 20--60\% of the domain height.  Episodes terminate at the
lower or upper altitude bound (10\% or 90\% of the domain height) or when the
wind sequence ends.  We do not define a 100 m success threshold; results are
reported as continuous distributions and, in the completed dynamic study,
will also include termination reasons.

\subsection{Quasi-steady discrete regime}
\label{sec:steady_environment}

The completed regime represents a glider with a quasi-steady polar.  For an
angle of attack $\alpha$ and bank angle $\mu$, the model obtains lift and drag
coefficients from the polar and solves for trim airspeed $V$, glide angle
$\gamma$, and turn rate.  A control change incurs a multiplicative drag penalty
of 1.1.  The controller acts every 1 s, with two physics integrations per
control step.  Wind data advance once every two control steps.  Thus the flight
model is quasi-steady, but the agent still flies through a time-varying DNS
sequence.

Actions are the nine combinations of incremental changes
$(\Delta\alpha,\Delta\mu)\in\{-1,0,1\}^{2}$.  Angle of attack spans
0--9 degrees in 3-degree increments, and bank spans $-20$ to 20 degrees in
10-degree increments.  The observation contains the current control indices,
the local vertical-wind acceleration $a_w$, and the difference in vertical wind
between the two wing tips $\Delta w$.  The reward is
\begin{equation}
  r_t = w_z(t) + 5a_w(t),
\end{equation}
where $w_z$ is local vertical wind.  This reward guides learning; it is not the
primary reported outcome.

Tabular Q-learning receives three hysteretic bins for each sensor.  DQN receives
the four-component continuous observation after standardizing the sensor
components with statistics collected from 20 random episodes.  The tabular
controller uses 100,000 interaction steps, $\gamma=0.98$, and an
$\epsilon$-greedy schedule from 1.0 to 0.01.  The selected DQN uses 200,000
interaction steps, learning rate $10^{-4}$, $\gamma=0.995$, a replay buffer of
100,000 transitions, batch size 128, target-network update frequency 2,000,
and an exploration fraction of 0.3.  Shared runtime hyperparameters are defined
only in \texttt{config.py}.

\subsection{Dynamic aerodynamic regime}
\label{sec:dynamic_environment}

The dynamic regime tests the same local energy-harvesting question after
removing the quasi-steady lookup approximation.  Its state comprises three-
dimensional position, ground velocity, angle of attack, bank angle, and a
continuous wind-frame coordinate.  At each 0.1 s integration substep, the
environment obtains temporally interpolated wind, forms air-relative velocity,
and computes lift, drag, and gravity from the same polar used in the discrete
regime.  A fourth-order Runge--Kutta integrator advances the coupled flight and
wind-frame state.  This design preserves the DNS input and wind-scaling rule
while exposing the controller to acceleration, transient airspeed, and
time-varying aerodynamics absent from the quasi-steady model.

The dynamic observation is
\begin{equation}
  o_t=[H_E,\dot H_E^{\mathrm{LP}},\Delta w_n,\mu],
\end{equation}
where $H_E=z+V_{\mathrm{TAS}}^2/(2g)$ is total-energy height,
$\dot H_E^{\mathrm{LP}}$ is a one-second low-pass estimate of its rate of
change, $\Delta w_n$ is the wingtip wind difference projected onto the lift
normal, and $\mu$ is bank angle.  The reward is the one-step increment in
$H_E$.  This choice makes the training signal and primary energy outcome
directly aligned while retaining altitude change as a complementary, more
interpretable measure.

The native action has two bounded commands: angle of attack and bank.  Both
commands pass through first-order actuators with separate time constants and
rate limits before affecting the vehicle.  PPO receives this continuous action
interface through a squashed Gaussian policy.  Dynamic DQN instead operates on
a $3\times3$ grid of angle-of-attack and bank commands.  The two learned
policies therefore share dynamics, observations, reward, initial-scenario
sampling, and wind data, but not the same action parameterization.  Their
comparison is evidence about controller variants in this environment, not an
action-space-controlled claim that one RL algorithm is intrinsically better.

Episodes terminate on lower or upper altitude boundaries, numerical divergence,
or exhaustion of the DNS sequence.  A low true airspeed disables the
aerodynamic-force calculation for that integration substep but does not end the
episode; numerical divergence is explicitly recorded.  The final dynamic
analysis will report each termination category.  In particular, termination at
the upper altitude boundary indicates that the simulated episode ended after
energy gain, but does not establish sustained soaring or a safety property.

\section{Experimental protocol}
\label{sec:protocol}

The steady study compares random control, tabular Q-learning, and DQN on 100
matched scenarios.  A scenario fixes the start frame, initial position,
altitude, and heading; all policies in that row receive the same scenario.
Scenario generation uses evaluation seed 20260816, and the random-policy action
stream is separately seeded.  The reported outcome is final minus initial
height in metres.  We report the mean, median, sample standard deviation, and
percentile bootstrap 95\% confidence interval for the mean using 20,000
resamples.  Matched-scenario differences are summarized with the same bootstrap
procedure.

The released steady artifacts contain one selected training run per trainable
policy under training seed 1.  Consequently, uncertainty intervals describe the
finite matched-scenario sample, not variation over training runs.  We reserve
multi-seed retraining and a separately generated test-scenario suite as a
mandatory follow-up before making confirmatory comparative claims.

\subsection{DQN-informed sensor bins}

After training the selected DQN, we sampled its action map on a standardized
$21\times21$ grid over $(a_w,\Delta w)$ and fit two thresholds per sensor to
its preferred action components.  The resulting raw-sensor boundaries are
$[-0.8,0.25]$ for $a_w$ and $[-0.15,0.60]$ for $\Delta w$; the action-map fit
error was 0.088.  These thresholds replace the original tabular bins, while
the tabular training budget and evaluation protocol remain unchanged.  This is
a representation-design experiment, not policy distillation: the tabular agent
is trained independently from a fresh Q table.

The procedure is exploratory because the tuned DQN and the derived bins were
selected using the same 100 scenarios subsequently summarized below.  The
derived-bin result is included to document the hypothesis and to guide the
next experiment, not as an out-of-sample comparison.

\subsection{Dynamic-regime setup}
\label{sec:dynamic_protocol}

Dynamic PPO and DQN use the four-dimensional observation above, with each
dimension standardized by random-rollout statistics stored with the resulting
model artifact.  PPO is trained for 100,000 transitions in a 64-environment
batch, using 16-step rollouts, GAE, clipped policy updates, and an entropy
coefficient annealed from 0.01 to zero.  Dynamic DQN is trained for 150,000
transitions with a 100,000-transition replay buffer, batch size 64, Double-DQN
targets, Huber loss, and a 3-by-3 action grid.  The matched evaluation includes
random-grid and fixed-command cruise controls as baselines, together with PPO
and DQN.  For every scenario, we retain height change, total-energy-height
change, episode length, and termination reason.

The existing 100-scenario dynamic result is an exploratory checkpoint.  The
completion experiment trained three independent replicas of each policy with
seeds 11, 22, and 33, then evaluated all frozen replicas on a new scenario
suite generated with evaluation seed 20260821.  This separates model selection
from the paper's comparative estimates.  Aggregation treats training seed as
the independent replication unit and reports scenario distributions within each
seed; it does not pool all episodes across seeds as though they were independent
training runs.

\section{Results: quasi-steady discrete regime}
\label{sec:steady_results}

The original matched-scenario comparison shows a modest ordering in favor of
the continuous-observation DQN.  Its mean height change was 42.4 m, whereas
random control lost 107.1 m on average (Table~\ref{tab:steady_baseline}).  The
paired DQN-minus-random difference was 149.6 m (bootstrap 95\% CI
$[74.2,224.4]$).  The corresponding tabular-Q-minus-random difference was
102.0 m (95\% CI $[25.7,178.3]$).  The DQN mean confidence interval still
crosses zero, and all three policies show substantial scenario-to-scenario
variation.  This baseline therefore supports an energy-harvesting signal in the
matched sample but does not establish reliable performance across training
seeds.

\begin{table}[t]
  \caption{Original steady-regime evaluation on 100 matched scenarios. Height
  change is final minus initial altitude in metres. Intervals are percentile
  bootstrap 95\% confidence intervals for the mean.}
  \label{tab:steady_baseline}
  \centering
  \small
  \begin{tabular}{lrrrr}
    \toprule
    Policy & Mean [95\% CI] & Median & SD & $n$ \\
    \midrule
    Random & $-107.1$ [$-162.8,-48.7$] & $-186.2$ & $292.8$ & 100 \\
    Tabular Q & $-5.1$ [$-72.0,64.0$] & $-148.4$ & $348.7$ & 100 \\
    DQN & $42.4$ [$-24.8,111.4$] & $49.2$ & $349.4$ & 100 \\
    \bottomrule
  \end{tabular}
\end{table}

Hyperparameter tuning substantially improved the selected DQN on the same
scenario set, yielding a mean height change of 337.4 m (95\% CI
$[297.4,376.2]$).  The DQN-informed discretization also improved the tabular
result to 192.3 m (95\% CI $[130.0,252.5]$; Table~\ref{tab:exploratory}).  The
median of the derived-bin tabular controller, 280.5 m, makes the improvement
visible beyond the mean.  However, these numbers must not be read as a
confirmatory superiority result.  The process selected both the DQN candidate
and the sensor thresholds using this evaluation suite, and it used a single
training seed.  The appropriate interpretation is that DQN action boundaries
are a promising hypothesis for designing a compact tabular state space.

\begin{table}[t]
  \caption{Exploratory steady-regime results on the same 100 scenarios used
  during selection. These results are not held-out evidence.}
  \label{tab:exploratory}
  \centering
  \small
  \begin{tabular}{lrrrr}
    \toprule
    Policy & Mean [95\% CI] & Median & SD & $n$ \\
    \midrule
    Tuned DQN & $337.4$ [$297.4,376.2$] & $352.1$ & $201.0$ & 100 \\
    DQN-informed tabular Q & $192.3$ [$130.0,252.5$] & $280.5$ & $315.5$ & 100 \\
    \bottomrule
  \end{tabular}
\end{table}

\section{Results: dynamic aerodynamic regime}
\label{sec:dynamic_results}

The dynamic evaluation used 100 matched scenarios generated with evaluation
seed 20260821, which was not used to select a model or tune a hyperparameter.
Each learned controller was trained independently with seeds 11, 22, and 33;
all six frozen artifacts were then evaluated sequentially against the same
scenario rows.  The replication unit is the training seed, while the 100
scenarios quantify within-replica variation.  Table~\ref{tab:dynamic_results}
reports the mean, median, sample standard deviation, and scenario-bootstrap
interval for each replica.  Figure~\ref{fig:dynamic_results} shows the full
matched-scenario distributions.

PPO produced positive mean height and total-energy-height changes for all three
training seeds.  Its mean height changes were $397.7$, $276.8$, and $378.5$ m,
with a mean of seed means of $351.0$ m and a between-seed standard deviation of
$65.0$ m.  The corresponding total-energy-height changes were $396.5$,
$278.1$, and $378.6$ m.  PPO therefore preserved energy gain after replacing
the quasi-steady lookup with integrated flight and actuator dynamics in this
held-out suite.

Dynamic DQN was less consistent in absolute energy gain.  Its three mean
height changes were $-146.3$, $-121.6$, and $-76.0$ m, and its mean of seed
means was $-114.6$ m (between-seed SD $35.7$ m).  Its mean total-energy-height
change was $-98.2$ m (between-seed SD $36.3$ m).  Its three-seed average was
less negative than the Random-grid baseline ($-125.1$ m height, $-117.6$ m
total-energy height) and the fixed Cruise baseline ($-167.2$ m height,
$-165.6$ m total-energy height),
but did not yield net positive energy-height change.  Because PPO uses a
continuous two-command action and DQN uses a $3\times3$ discrete grid, this is
a comparison of controller variants with different action interfaces, not a
pure algorithm ranking.

Termination patterns provide important context.  PPO reached the upper
altitude boundary in 45--64\% of scenarios across seeds, with wind-end
fractions of 31--46\%.  DQN reached the lower altitude boundary in 50--65\% of
scenarios and the upper boundary in only 4--5\%; numerical-divergence fractions
were at most 1\%.  Random grid and Cruise reached the lower boundary in 66\%
and 73\% of scenarios, respectively.  An upper-bound termination records that
an episode ended after gaining energy in the finite domain; it is not evidence
of sustained soaring, safety, or indefinite endurance.

\begin{table}[t]
  \caption{Held-out dynamic evaluation on 100 matched scenarios. Each learned
  row is one independent training seed; intervals are percentile bootstrap 95\%
  CIs for the within-seed scenario mean. Height and total-energy-height changes
  are in metres.}
  \label{tab:dynamic_results}
  \centering
  \small
  \begin{tabular}{llrrrr}
    \toprule
    Policy & Seed & Height mean [95\% CI] & Energy mean & Median & SD \\
    \midrule
    Random grid & -- & $-125.1$ [$-186.3,-59.9$] & $-117.6$ & $-222.3$ & $322.7$ \\
    Cruise & -- & $-167.2$ [$-216.6,-114.0$] & $-165.6$ & $-238.5$ & $261.1$ \\
    DQN & 11 & $-146.3$ [$-197.4,-92.6$] & $-131.1$ & $-208.3$ & $265.9$ \\
    DQN & 22 & $-121.6$ [$-173.1,-67.9$] & $-104.3$ & $-198.4$ & $272.1$ \\
    DQN & 33 & $-76.0$ [$-132.5,-17.3$] & $-59.3$ & $-173.7$ & $297.7$ \\
    PPO & 11 & $397.7$ [$351.3,440.2$] & $396.5$ & $442.8$ & $229.1$ \\
    PPO & 22 & $276.8$ [$212.9,338.6$] & $278.1$ & $399.7$ & $326.7$ \\
    PPO & 33 & $378.5$ [$335.4,418.7$] & $378.6$ & $417.8$ & $212.8$ \\
    \bottomrule
  \end{tabular}
\end{table}

\begin{figure}[t]
  \centering
  \includegraphics[width=\linewidth]{../trainresult/dynamic_multiseed/dynamic_multiseed_results.png}
  \caption{Matched-scenario distributions for the held-out dynamic evaluation.
  Points retain scenario-level variation; learned-policy points include all
  three training-seed replicas.}
  \label{fig:dynamic_results}
\end{figure}

\section{Discussion and limitations}
\label{sec:discussion}

The two fidelity regimes form a progression rather than a cross-regime
leaderboard.  The quasi-steady environment isolates local energy-harvesting
decisions with transparent discrete dynamics and was useful for exposing a
DQN-informed tabular representation.  The dynamic environment then tests the
same DNS-driven task with continuous-time flight integration, temporal wind
interpolation, actuator lag, and a reward defined directly by total-energy
height.  The held-out dynamic results show that PPO's positive energy-height
signal survives this increase in model fidelity across all three training
seeds.  Dynamic DQN remained negative in absolute energy-height change, despite
a less-negative three-seed average than the non-learning baselines, suggesting
that the coarse action grid or value-learning configuration is not yet
sufficient for the richer control problem.  Since the action interfaces differ,
this observation should
not be interpreted as an intrinsic PPO-versus-DQN result.

Several limitations bound the claim.  The DNS flow is a controlled numerical
input at $Ra=5\times10^7$, scaled to the simulated vehicle; it is not an
atmospheric measurement.  The glider remains a point-mass model driven by a
fixed aerodynamic polar, and the dynamic observation contains only local
quantities rather than the full flow field.  We used three training seeds and
one held-out 100-scenario suite; the bootstrap intervals quantify scenario
sampling, whereas the small between-seed sample limits inference about training
variation.  Upper-altitude termination is common for PPO and truncation at the
end of the finite wind sequence is also common, so the results demonstrate
energy gain over finite episodes rather than sustained safe soaring.  No route
objective, collision model, deployment constraint, real-flight experiment, or
sim-to-real transfer is evaluated.

\section{Conclusion}
\label{sec:conclusion}

ThermalHunter provides a two-fidelity, DNS-driven platform for studying
energy-harvesting soaring in time-varying turbulent wind.  In the completed
held-out dynamic experiment, PPO replicas achieved positive height and
total-energy-height changes across all three training seeds, while dynamic DQN
improved over Random-grid and Cruise controls without achieving positive mean
energy gain.  The results support the narrower claim that reinforcement
learning can discover energy-harvesting behavior in this controlled numerical
environment.  They do not establish route navigation, safety, deployment
readiness, sustained soaring, or transfer beyond the simulated flow.  Future
work should expand training-seed replication, vary the held-out DNS scenarios,
and test action representations and dynamic models under explicitly matched
control interfaces.

\begin{thebibliography}{9}

\bibitem[Colabrese et al.(2017)]{colabrese2017}
S. Colabrese, K. Gustavsson, A. Celani, and L. Biferale.
Flow navigation by smart microswimmers via reinforcement learning.
\emph{Physical Review Letters}, 118:158004, 2017.
doi: \url{https://doi.org/10.1103/PhysRevLett.118.158004}.

\bibitem[Mnih et al.(2015)]{mnih2015}
V. Mnih, K. Kavukcuoglu, D. Silver, A. A. Rusu, J. Veness, M. G. Bellemare,
A. Graves, M. Riedmiller, A. K. Fidjeland, G. Ostrovski, et al.
Human-level control through deep reinforcement learning.
\emph{Nature}, 518:529--533, 2015.
doi: \url{https://doi.org/10.1038/nature14236}.

\bibitem[Novati et al.(2019)]{novati2019}
G. Novati, H. Mahadevan, and P. Koumoutsakos.
Controlled gliding and perching through deep-reinforcement-learning.
\emph{Physical Review Fluids}, 4:093902, 2019.
doi: \url{https://doi.org/10.1103/PhysRevFluids.4.093902}.

\bibitem[Reddy et al.(2018)]{reddy2018}
G. Reddy, A. F. E. Schaefer, M. D. T. P. de~Wilde, C. A. V. de~Gelder,
and I. D. Couzin.
Glider soaring via reinforcement learning in the field.
\emph{Nature}, 562:236--239, 2018.
doi: \url{https://doi.org/10.1038/s41586-018-0533-0}.

\end{thebibliography}

\end{document}
